% Este arquivo é gerado por scripts/generate_resume.rb.
% Edite curriculo-geral.yaml, não este arquivo.

\resumeheader{Marcio Valente}{Tecnologia da Informação | Sistemas | Suporte Técnico}{\href{\detokenize{mailto:eric.melv42@gmail.com}}{eric.melv42@gmail.com}}{(92) 99514-4499}{Manaus, Amazonas, Brasil}{\href{\detokenize{https://www.linkedin.com/in/marcioericlvalente}}{linkedin.com/in/marcioericlvalente}}{\href{\detokenize{https://www.github.com/mericxy}}{github.com/mericxy}}

\sectiontitle{Resumo Profissional}
\resumeparagraph{Profissional de tecnologia da informação e estudante de Engenharia de Software, com experiência em desenvolvimento de sistemas web, apoio técnico e administrativo, organização de demandas, documentação e resolução de problemas. Atuo em projeto de pesquisa da Motorola em parceria com a UFAM, colaborando na criação e manutenção de ferramentas, dashboards e sistemas internos. Busco oportunidades em suporte, sistemas, operações de TI ou desenvolvimento, com facilidade para aprender ferramentas e trabalhar em equipe.}

\sectiontitle{Competências Técnicas}
\skill{Suporte e operações}{apoio técnico, resolução de problemas, organização de rotinas, documentação e atendimento de demandas.}
\skill{Sistemas e dados}{aplicações web, integração de APIs, dashboards, visualização de dados, PostgreSQL e deploy em servidor.}
\skill{Desenvolvimento}{React, JavaScript, TypeScript, Python, Django, FastAPI, Node.js, HTML5 e CSS3.}
\skill{Ferramentas e processos}{Git, GitHub, Docker, Jira, Figma, testes de software, trabalho em equipe ágil e entregas incrementais.}

\sectiontitle{Experiência Profissional}
\role{Frontend Developer}{Mai/2025 -- Presente}{MOB4AI | Projeto de pesquisa Motorola/UFAM}{Manaus, Amazonas}
\achievement{Desenvolvi interfaces interativas com React para visualização de dados reais em projeto de pesquisa em parceria entre Motorola e UFAM.}
\achievement{Atuei na melhoria da experiência de uso em dashboards analíticos, com foco em navegação, legibilidade e interação entre visualizações de dados.}
\achievement{Construí funcionalidades orientadas a dados usando bibliotecas de visualização como Apache ECharts e Highcharts.}
\achievement{Participei de decisões de arquitetura frontend, definição de stack para novas ferramentas e deploy em servidor em contato com outros times.}
\vspace{4pt}
\role{Estagiário}{Nov/2024 -- Mai/2025}{CETAM}{Manaus, Amazonas}
\achievement{Prestei apoio técnico e administrativo em ambiente corporativo, desenvolvendo organização, comunicação, resolução de problemas operacionais e adaptação a rotinas profissionais.}

\sectiontitle{Projetos}
\project{DocParse | \href{\detokenize{https://docparse.meric.dev.br/}}{docparse.meric.dev.br} | \href{\detokenize{https://github.com/mericxy/DocParse}}{github.com/mericxy/DocParse}}{Processamento e transcrição de documentos | Full Stack e OCR}{Aplicação web para leitura, OCR e extração estruturada de dados de cartões de ponto e holerites em PDF para planilhas. Inclui processamento assíncrono em Python, interface de conferência em React e deploy conteinerizado com Docker.}

\project{ReservaUFAM | \href{\detokenize{https://github.com/mericxy/ReservaUFAM}}{github.com/mericxy/ReservaUFAM}}{Sistema de reservas institucionais | Projeto em equipe}{Projeto desenvolvido para disciplinas da UFAM com o objetivo de centralizar reservas de recursos institucionais. Atuei na construção das interfaces e na integração das funcionalidades do sistema.}

\project{NAVIR Web | \href{\detokenize{https://github.com/navir-ufam/NAVIR-Web}}{github.com/navir-ufam/NAVIR-Web}}{Sistema de gestão acadêmica | Liderança técnica}{Sistema interno do laboratório NAVIR da UFAM. Como tech lead, colidero a frente técnica, apoiando decisões de arquitetura, organização de demandas e desenvolvimento de módulos para gestão acadêmica, controle de acesso, dashboards e relatórios.}

\sectiontitle{Educação}
\role{Bacharelado em Engenharia de Software}{Out/2022 -- Conclusão prevista: Jan/2027}{Universidade Federal do Amazonas}{Manaus, Amazonas}

\sectiontitle{Idiomas}
\resumelanguage{Português}{Nativo}
\resumelanguage{Inglês}{Intermediário para leitura de documentação técnica, escrita no dia a dia e comunicação funcional}

\sectiontitle{Informações Adicionais}
\resumeparagraph{Carteira Nacional de Habilitação (CNH): categorias A e B}
